\documentclass{article} 
\input{../../../lib.sty} 
 
\title{\texttt{fn sample\_geometric\_exp\_fast}} 
\author{\MichaelShoemateAuthor} 
 
\begin{document} 
\maketitle 
 
Proves soundness of \texttt{fn sample\_geometric\_exp\_fast} in \asOfCommit{mod.rs}{0be3ab3e6}. 
This proof is an adaptation of \href{https://arxiv.org/pdf/2004.00010.pdf#subsection.5.2}{subsection 5.2} of \cite{CKS20}. 
 
\section{Hoare Triple} 
 
\subsection*{Precondition} 
\subsubsection*{Compiler-verified} 
\begin{itemize}
    \item Argument \texttt{x} is of type \texttt{RBig}, a bignum rational.
\end{itemize}
 
\subsubsection*{Caller-verified} 
\begin{itemize}
    \item $\texttt{x} \geq 0$
\end{itemize}
 
\subsection*{Pseudocode} 
\lstinputlisting[language=Python,firstline=2,escapechar=|]{./pseudocode/sample_geometric_exp_fast.py} 
 
\subsection*{Postcondition} 
\begin{theorem}
\label{postcondition} 
For any setting of the input parameter \texttt{x} such that the given preconditions hold, \\ 
\texttt{sample\_geometric\_exp\_fast} either returns \texttt{Err(e)} due to a lack of system entropy, 
or \texttt{Ok(out)}, where \texttt{out} is distributed as $\textrm{Geometric}(1 - \exp(-x))$. 
\end{theorem}

Assume the preconditions are met. 
 
\begin{lemma}\label{x-parts} 
    There exist a non-negative integer $s$ and positive integer $t$ such that $x = s/t$. 
    If $x > 0$, then $s \geq 1$. 
\end{lemma} 
 
\begin{proof} 
    Since \texttt{x} is of type \docsrs{dashu}{rational/struct}{Rbig} and $\texttt{x} \geq 0$, 
    it can be written as $x = s/t$ for some non-negative integer $s$ and positive integer $t$. 
    This is why \texttt{RBig.into\_parts} is infallible. 
    If $x > 0$, then its numerator is positive, so $s \geq 1$. 
\end{proof} 
 
\begin{lemma}\label{err-e} 
    \texttt{sample\_geometric\_exp\_fast} only returns \texttt{Err(e)} when there is a lack of system entropy. 
\end{lemma} 
 
\begin{proof} 
    If $x = 0$, the function returns \texttt{Ok(0)} immediately, so it cannot return \texttt{Err(e)}. 
    Assume $x > 0$. By Lemma \ref{x-parts}, there exist integers $s \geq 1$ and $t > 0$ with $x = s/t$. 
    Since $t$ is a positive integer, the preconditions on \rustdoc{traits/samplers/uniform/fn}{sample\_uniform\_ubig\_below} are met. 
    Any returned value \texttt{u} therefore lies in $\{0, 1, \dots, t - 1\}$, so $u/t \geq 0$, 
    and the preconditions on \rustdoc{traits/samplers/cks20/fn}{sample\_bernoulli\_exp} are also met. 
    The input on line \ref{line:V} is the non-negative rational $1$, so the preconditions on 
    \rustdoc{traits/samplers/cks20/fn}{sample\_geometric\_exp\_slow} are met as well. 
    By their definitions, each of these fallible calls can only return an error due to a lack of system entropy. 
    These are the only fallible calls in \texttt{sample\_geometric\_exp\_fast}, 
    so the function only returns \texttt{Err(e)} when there is a lack of system entropy. 
\end{proof} 
 
We now establish some lemmas that will be useful in proving the distribution of \texttt{out}. 
 
\begin{lemma}\label{loop-dist} 
    Suppose $x > 0$. Conditioned on the loop on lines \ref{line:U}--\ref{line:D} terminating without returning \texttt{Err(e)}, 
    the value \texttt{u} at loop exit is a realization of a random variable $U_{\mathrm{acc}}$ supported on $\{0, 1, \dots, t - 1\}$, where 
    \[
        P[U_{\mathrm{acc}} = u] = \frac{1 - e^{-1/t}}{1 - e^{-1}} e^{-u/t}. 
    \]
\end{lemma} 
 
\begin{proof} 
    By Lemma \ref{x-parts}, $t > 0$, so in each loop iteration the call on line \ref{line:U} samples 
    a random variable $U \sim \textrm{Uniform}(0, t)$ supported on $\{0, 1, \dots, t - 1\}$. 
    Conditioned on a realization $u$ of $U$, the call on line \ref{line:D} samples 
    $D \sim \textrm{Bernoulli}(\exp(-u/t))$ by the specification of 
    \rustdoc{traits/samplers/cks20/fn}{sample\_bernoulli\_exp}. 
 
    Each iteration uses fresh randomness, and the loop stops at the first iteration for which $D = \top$. 
    Therefore the value returned by the loop has the same distribution as $U$ conditioned on $D = \top$. 
    For any $u \in \{0, 1, \dots, t - 1\}$, 
    \begin{align*} 
        P[U_{\mathrm{acc}} = u] 
        &= P[U = u \mid D = \top] \\ 
        &= \frac{P[U = u] P[D = \top \mid U = u]}{P[D = \top]} && \text{by Bayes' theorem} \\ 
        &= \frac{(1/t)e^{-u/t}}{\frac{1}{t} \sum_{k=0}^{t-1} e^{-k/t}} \\ 
        &= \frac{(1/t)e^{-u/t}}{\frac{1}{t}\frac{1 - e^{-1}}{1 - e^{-1/t}}} \\ 
        &= \frac{1 - e^{-1/t}}{1 - e^{-1}} e^{-u/t}. 
    \end{align*} 
\end{proof} 
 
\begin{lemma}\label{z-dist} 
    Suppose $x > 0$ and \texttt{sample\_geometric\_exp\_fast} returns \texttt{Ok(out)}. 
    Then the value \texttt{z} is a realization of a random variable $Z \sim \textrm{Geometric}(1 - \exp(-1/t))$. 
\end{lemma} 
 
\begin{proof} 
    By Lemma \ref{loop-dist}, the value \texttt{u} exiting the loop is distributed as $U_{\mathrm{acc}}$. 
    The call on line \ref{line:V} returns a value \texttt{v} distributed as 
    $V \sim \textrm{Geometric}(1 - \exp(-1))$ by the specification of 
    \rustdoc{traits/samplers/cks20/fn}{sample\_geometric\_exp\_slow}. 
    Since line \ref{line:V} is executed after the loop and uses fresh randomness, $U_{\mathrm{acc}}$ and $V$ are independent. 
 
    For any $z$, define $u_z := z \mod t$ and $v_z := \lfloor z/t \rfloor$, so that $z = u_z + t \cdot v_z$. 
    Then 
    \begin{align*} 
        P[Z = z] 
        &= P[U_{\mathrm{acc}} = u_z, V = v_z] && \text{since } z = u_z + t \cdot v_z \\ 
        &= P[U_{\mathrm{acc}} = u_z] P[V = v_z] && \text{by independence} \\ 
        &= \frac{1 - e^{-1/t}}{1 - e^{-1}} e^{-u_z/t} \cdot (1 - e^{-1}) e^{-v_z} && \text{by Lemma \ref{loop-dist}} \\ 
        &= (1 - e^{-1/t}) e^{-(u_z/t + v_z)} \\ 
        &= (1 - e^{-1/t}) e^{-z/t}. 
    \end{align*} 
    Therefore $Z \sim \textrm{Geometric}(1 - \exp(-1/t))$. 
\end{proof} 
 
\begin{lemma}\label{divide_geometric}\cite{CKS20} 
    Fix $p \in (0, 1]$. Let $G$ be a $\textrm{Geometric}(1 - p)$ random variable, and $n \geq 1$ be an integer. 
    Then $\lfloor G / n \rfloor$ is a $\textrm{Geometric}(1 - q)$ random variable with $q = p^n$. 
\end{lemma} 
 
\begin{proof} 
    \begin{align*} 
        P[\lfloor G/n \rfloor = k] &= P[nk \le G < (k + 1)n] &&\text{any $G$ in the interval maps to $k$} \\ 
        &= \sum_{l=kn}^{(k+1)n - 1} (1 - p)p^l \\ 
        &= (1 - p^n)p^{nk} \\ 
        &= (1 - q)q^k 
    \end{align*} 
\end{proof} 
 
\begin{lemma}\label{ok-out} 
    If the outcome of \texttt{sample\_geometric\_exp\_fast} is \texttt{Ok(out)}, 
    then \texttt{out} is distributed as $\textrm{Geometric}(1 - \exp(-x))$. 
\end{lemma} 
 
\begin{proof} 
    If $x = 0$, the function returns \texttt{Ok(0)} immediately. Since 
    $\textrm{Geometric}(1 - \exp(-0)) = \textrm{Geometric}(0)$ is the point mass at $0$, 
    the claim holds. 
 
    Assume $x > 0$. By Lemma \ref{x-parts}, there exist integers $s \geq 1$ and $t > 0$ with $x = s/t$. 
    By Lemma \ref{z-dist}, \texttt{z} is a realization of $Z \sim \textrm{Geometric}(1 - \exp(-1/t))$. 
    Applying Lemma \ref{divide_geometric} with $p = \exp(-1/t)$ and $n = s$, 
    we conclude that $\lfloor Z / s \rfloor$ is distributed as $\textrm{Geometric}(1 - \exp(-s/t))$. 
    Finally, $\texttt{out} = \lfloor Z / s \rfloor$ and $s/t = x$, so 
    \texttt{out} is distributed as $\textrm{Geometric}(1 - \exp(-x))$. 
\end{proof} 
 
\begin{proof} 
    Theorem~\ref{postcondition} holds by \ref{err-e} and \ref{ok-out}. 
\end{proof} 
 
\bibliographystyle{alpha} 
\bibliography{mod} 
\end{document}
